\documentclass[12pt,a4paper]{article}
\usepackage[french]{babel}
\usepackage[T1]{fontenc}
\usepackage{amsmath, amssymb}
\usepackage{geometry}
\geometry{top=1cm, bottom=2cm, left=1cm, right=1cm}
\usepackage{multicol}
\usepackage{graphicx}
\usepackage{tcolorbox}
\usepackage{multicol}

\begin{document}

\centering\textbf{Devoir Surveill\'e de Math\'ematiques - Statistiques et Probabilit\'es}\\
\vspace{0.5cm}

\noindent\textbf{Nom :} \underline{\hspace{5cm}} \hfill \textbf{Pr\'enom :} \underline{\hspace{5cm}} \\
\textbf{Classe :} \underline{\hspace{5cm}}\\
\vspace{1cm}

\centering\textit{\textbf{Chaque question vaut 1 point. La pr\'esentation et la rigueur seront \'egalement \'evalu\'ees. Calculatrice interdite.}}

\section*{Partie 1 : Automatismes}

\subsection*{Exercice 1 : QCM - Répondre sans justification}

\begin{enumerate}
    \item Un prix augmente de 20\% puis diminue de 20\%. Après ces deux évolutions, on peut affirmer que :\\
    a) Le prix est égal à sa valeur de départ \quad
    b) Le prix est strictement supérieur à sa valeur de départ \\
    c) Le prix est strictement inférieur à sa valeur de départ \quad
    d) Cela dépend de la valeur de départ

    \item Le tiers d’un quart correspond à la fraction : \\
    a) $\dfrac{1}{7}$ \quad b) $\dfrac{3}{4}$ \quad c) $\dfrac{1}{3} \times 4$ \quad d) $\dfrac{1}{12}$

    \item 50 élèves étudient le grec, ce qui représente 4\% des élèves du lycée. Le nombre total d'élèves est : \\
    a) 200 \quad b) 1250 \quad c) 125 \quad d) 2
\end{enumerate}

\vspace{0.5cm}
\begin{tcolorbox}[colframe=black, colback=white, width=\linewidth]
\vspace{6cm}
\end{tcolorbox}

\section*{Partie 2 : Fr\'equences marginales et conditionnelles}

\subsection*{Exercice 2 : Analyse d'un tableau de fr\'equences}

Dans une classe de 30 \`eleves, 12 sont des filles. Parmi elles, 8 aiment les math\'ematiques. Au total, 18 \`eleves aiment les math\'ematiques.

\begin{enumerate}
    \item Compl\'eter le tableau ci-dessous :
    
    \begin{center}
    \begin{tabular}{|c|c|c|c|}
    \hline
                 & Aiment les maths & N'aiment pas les maths & Total \\
    \hline
    Filles       &                   &                         &       \\
    \hline
    Gar\c{c}ons   &                   &                         &       \\
    \hline
    Total        &                   &                         & 30    \\
    \hline
    \end{tabular}
    \end{center}

    \item Quelle est la fr\'equence des \textbf{gar\c{c}ons} qui \textbf{aiment} les math\'ematiques ?
    \item Quelle est la fr\'equence des \textbf{\`eleves} qui \textbf{n'aiment pas} les math\'ematiques ?
    \item Quelle est la fr\'equence des \textbf{filles parmi ceux qui aiment} les math\'ematiques ?
\end{enumerate}

\vspace{0.5cm}
\begin{tcolorbox}[colframe=black, colback=white, width=\linewidth]
\vspace{10cm}
\end{tcolorbox}

\section*{Partie 3 : Probabilit\'es et probabilit\'es conditionnelles}

\subsection*{Exercice 3 : Probl\`eme de probabilit\'es}

Dans une bo\^ite, il y a 4 boules rouges, 3 boules vertes et 3 boules bleues. On tire une boule au hasard.

\begin{enumerate}
    \item Quelle est la probabilit\'e de tirer :
    \begin{itemize}
        \item une boule rouge ?
        \item une boule bleue ?
    \end{itemize}
    \item On sait que la boule tir\'ee n'est \textbf{pas rouge}. Quelle est la probabilit\'e qu'elle soit bleue ?
    \item On tire une boule, on la remet, puis on tire \`a nouveau une boule. Quelle est la probabilit\'e d'obtenir deux boules de m\^eme couleur ?
\end{enumerate}

\vspace{0.5cm}
\begin{tcolorbox}[colframe=black, colback=white, width=\linewidth]
\vspace{6cm}
\end{tcolorbox}

\end{document}